%!TEX root = ../template.tex
%%%%%%%%%%%%%%%%%%%%%%%%%%%%%%%%%%%%%%%%%%%%%%%%%%%%%%%%%%%%%%%%%%%%
%% abstrac-pt.tex
%% NOVA thesis document file
%%
%% Abstract in Portuguese
%%%%%%%%%%%%%%%%%%%%%%%%%%%%%%%%%%%%%%%%%%%%%%%%%%%%%%%%%%%%%%%%%%%%

\typeout{NT FILE abstrac-pt.tex}%
O aumento da frequência dos incêndios florestais, representam uma ameaça crescente 
para áreas residenciais, 
especialmente aquelas rodeadas por vegetação densa. A avaliação do 
risco de incêndio e a implementação de estratégias de mitigação são essenciais 
para a proteção de casas e comunidades. Este estudo explora o uso de técnicas de 
aprendizagem profunda para automatizar a classificação do risco de incêndio em 
habitações com base na vegetação envolvente.

Com base num conjunto de dados de mais de 2.000 casas isoladas em Portugal Continental, 
fornecido pela Associação para o Desenvolvimento da Aerodinâmica Industrial (ADAI), 
este trabalho utiliza imagens aéreas obtidas através da API de dados abertos da 
Direção-Geral do Território (DGT). Redes Neuronais Convolucionais (CNNs) foram 
usadas para desenvolver o modelo de classificação, avaliando arquiteturas 
como ResNet-50, DenseNet-161 e EfficientNet-B5. A análise compara modelos 
treinados com imagens RGB e aqueles que integram dados espectrais adicionais, 
como Infrared-Green-Red (IRG) e o Índice de Vegetação de Diferença Normalizada 
(NDVI), para avaliar o impacto na precisão da classificação.

Para melhorar o desempenho do modelo, aplicaram-se técnicas de augmentação 
e uma abordagem de aprendizagem em conjunto (ensemble learning), garantindo maior 
robustez na classificação de diferentes tipos de vegetação. A metodologia proposta 
visa tornar a avaliação de risco mais eficiente, reduzindo a necessidade de 
trabalho manual e permitindo uma identificação mais rápida de áreas vulneráveis. 
Este estudo contribui para o avanço da análise automatizada do risco de incêndios, 
apoiando estratégias preventivas e reforçando a resiliência das comunidades em zonas de risco.

% Palavras-chave do resumo em Português
% \begin{keywords}
% Palavra-chave 1, Palavra-chave 2, Palavra-chave 3, Palavra-chave 4
% \end{keywords}
\keywords{
  Aprendizagem Automática \and
  Aprendizagem Supervisionada \and
  Aprendizagem Profunda \and
  Redes Neurais Convolucionais \and
  Sensoriamento Remoto
}
% to add an extra black line
